\chapter{Distributing your app}
\label{ch:distrib}

\begin{aside}
You built it. Now ship it. Static binaries, package
managers, container images, source releases --- pick one or
all.
\end{aside}

\section{Static binary}

The simplest distribution. \code{-static} at link time;
one file users download and run. Size: 5--20 MB depending
on what's linked.

\section{Package managers}

\begin{itemize}[leftmargin=*]
  \item DEB (Debian, Ubuntu): \code{cpack -G DEB}.
  \item RPM (Fedora, RHEL): \code{cpack -G RPM}.
  \item Homebrew (macOS): a formula.
  \item Chocolatey (Windows): a nuspec.
  \item AUR (Arch): a PKGBUILD.
  \item Nix: a derivation.
\end{itemize}

Each has its own packaging discipline. \maya{}'s built-in
CMake/CPack handles the first two; the rest are
user-written.

\section{Container image}

\begin{lstlisting}
FROM debian:bookworm-slim
COPY myapp /usr/local/bin/
ENTRYPOINT ["myapp"]
\end{lstlisting}

A 50 MB image; fine for CI or for distributing as
\code{docker run}.

\section{Source release}

A tarball with source and \code{CMakeLists.txt}. Users
build themselves. Smallest download, biggest user effort.

\section{Versioning}

Semver: \code{major.minor.patch}. Tag releases in git.
Embed the version in the binary via a \code{config.h}
generated from CMake.

\section{Release notes}

CHANGELOG.md with sections per version. Categories:
Added, Changed, Fixed, Removed.

\section{Signing}

GPG-sign release tarballs. Provide checksums (sha256).
Users verify before installing.

\section{Auto-update}

If your app self-updates, the channel is a security
boundary. Sign updates; verify on download.

\section{Recap}

\begin{recap}
\begin{itemize}[leftmargin=*]
  \item Static binary is simplest.
  \item Package managers reach more users.
  \item Containers for CI.
  \item Semver + changelog + signing.
\end{itemize}
\end{recap}

\section*{Workshop}

\begin{enumerate}[leftmargin=*]
  \item Build a static binary; check the size.
  \item Generate a DEB via CPack.
  \item Write a 5-line changelog for your last 5 commits.
\end{enumerate}
